% Cardiac arrest early warning pipeline: leakage prevention annex
% Standalone document answering the reviewer point on leakage prevention:
% whether the model recognizes patients instead of deterioration, identity
% invariance tests, and unseen hospital evaluation. Research prototype only:
% nothing here is clinical evidence. Style note: no hyphens or dashes in prose.

\documentclass[11pt]{article}

\usepackage[utf8]{inputenc}
\usepackage[T1]{fontenc}
\usepackage{amsmath}
\usepackage{graphicx}
\usepackage{booktabs}
\usepackage{caption}
\usepackage[hidelinks]{hyperref}
\usepackage[a4paper,margin=1in]{geometry}

\graphicspath{{images/}}

\title{Leakage Prevention: Identity Recognition, Invariance Tests and
Unseen Hospital Evaluation}
\author{[Author names to be added]}
\date{23 September 2026}

\begin{document}

\maketitle

\begin{center}
\small
Research prototype. All results are pipeline demonstrations, not clinical
validation. Identity probe accuracies are means over five cross validation
folds; chance is 0.050 (one in twenty training records).
\end{center}

\section{Does the model recognize patients instead of deterioration?}
We test this directly with a linear identity probe. For each cross validation
fold, the trained model encodes every training window into its context vector,
a logistic regression is fitted on 80 percent of those windows to classify the
patient, and the held out 20 percent are scored. A model whose representation
is invariant to patient identity would score at chance. Table~\ref{tab:probe}
reports the probe for all eight trained variants.

\begin{table}[h]
\centering
\caption{Linear identity probe on the context vector (chance 0.050).}
\label{tab:probe}
\begin{tabular}{lc}
\toprule
Variant & Probe accuracy \\
\midrule
FEAN, concatenation (baseline) & 0.423 \\
FEAN, identity adversarial (gradient reversal) & 0.470 \\
FEAN, identity adversarial + PCGrad & 0.438 \\
FEAN, gated fusion + masks & 0.402 \\
FEAN, static modality weights + masks & 0.354 \\
FEAN, modality attention + masks & 0.143 \\
Ablation: engineered features only & 0.130 \\
Ablation: embeddings only & 0.439 \\
\bottomrule
\end{tabular}
\end{table}

The answer is yes, with two qualifications. The models do recognize patients:
every embedding based variant scores well above chance, so leakage is present
and quantified rather than assumed. But the identity signal is carried almost
entirely by the embedding block: the probe falls from 0.439 to 0.130 when the
ECG FM embeddings are replaced by engineered features alone. Modality
attention is the only architecture that substantially reduces leakage (0.143),
and it is also the best held out variant at 6 hours (AUROC 0.473). The two
observations agree: identity laden embeddings hurt generalization, and
mechanisms that down weight them help.

\section{Identity invariance tests}
Two forms of invariance testing are in place.

Measurement. The linear probe above is the invariance test, computed per fold
for every variant, so invariance is a measured quantity, not an assumption.

Enforcement. Identity adversarial training (a gradient reversal layer against
a patient ID classifier, optionally with PCGrad across task gradients) was
added to enforce invariance. The honest result: it improves held out AUROC
slightly at both horizons but does not reduce the probe (0.423 to 0.470), so
gradient reversal did not remove the leakage on this data. This negative
result is reported in the manuscript rather than hidden. A natural next step,
not yet run, is to combine the two mechanisms that do work: modality
attention, which down weights the identity carrying embeddings, together with
an identity adversary applied to the embedding block only.

\section{Unseen hospital evaluation}
What is covered today:

\begin{itemize}
\item Unseen patients: all training and evaluation splits are grouped by
patient, and a patient never appears in more than one split (asserted in
code).
\item Unseen domain: the synthetic trained model is applied zero shot to the
real Sudden Cardiac Death Holter Database and scores at chance (1 hour AUROC
0.338, 6 hour AUROC 0.447), which quantifies the synthetic to real gap.
\end{itemize}

What is not covered: an evaluation on a second hospital. The Sudden Cardiac
Death Holter Database is the only open access dataset with hours of pre arrest
recording, so no second open site exists for this task, and the test cannot be
run on open data alone.

The planned unseen hospital evaluation is the credentialed MIMIC III Ext CA
set: 36 arrest episodes in 31 ICU patients with photoplethysmogram plus
electrocardiogram or arterial pressure, a different site, era, monitor type
and modality set from the Holter database. The loader is written and self
tested (window manifest with resuscitation exclusions, patient level splits,
loud failure until credentialed access lands); the evaluation itself is
pending access approval and will be reported with the same patient level,
confidence interval aware protocol.

\section{Summary}
Patient recognition is measured, not assumed: the models recognize patients,
the identity signal lives in the embedding block, gradient reversal does not
remove it, and modality attention does. Unseen patient and unseen domain
evaluation exist; unseen hospital evaluation is designed, loader ready, and
blocked only on credentialed data access.

\end{document}